\documentclass[11pt,letterpaper]{article}
\usepackage{amsmath,amssymb,mathtools}
\usepackage[hidelinks]{hyperref}
\hypersetup{pdftitle={F0: National reference program},
  pdfauthor={Agent: modeler (openai-codex/gpt-6-astra)}}
\setlength{\textwidth}{6.5in}
\setlength{\oddsidemargin}{0in}
\setlength{\evensidemargin}{0in}
\setlength{\textheight}{9in}
\setlength{\topmargin}{-0.5in}
\setlength{\parskip}{0.35em}
\setlength{\parindent}{0pt}
\renewcommand{\theequation}{F0.\arabic{equation}}
\newcommand{\issue}[1]{\href{https://github.com/helios1168/td/issues/#1}{\##1}}
\newcommand{\MODEL}{\href{MODEL.md}{\texttt{docs/MODEL.md}}}
\DeclareMathOperator{\capcut}{cap}
\DeclareMathOperator{\lexmin}{lexmin}
\title{F0: National reference program\\\large Authoritative mathematical statement}
\author{Agent: modeler (openai-codex/gpt-6-astra)}
\date{15 October 2026}
\begin{document}
\maketitle

\textbf{Domain and status.} F0 is the national (main-map) reference problem,
not an IFA formulation or an algorithm. This file is the single authoritative
statement of F0; model-sync edits this source and rebuilds \texttt{F0.pdf}.
It replaces the mathematics formerly in \MODEL, \S0, preserving that program
and incorporating the owner's 15 October Manhattan ruling on \issue{140}.
Its relation to F0 is \emph{exact}, by definition. A feasibility certificate
covers the complete routed national ledger and every hard condition below,
including the explicitly identified Manhattan exception. An optimality
certificate additionally needs numerical owner-defined shape and balance
evaluators and all three lexicographic stages. The production pipeline supplies
neither a national F0 optimizer nor such an optimality certificate.

\section{Sets, data and notation}
The ground set $Z$ is the explicitly shipped set of 2025 CONUS ZCTAs,
including zero-opportunity territory. One multipart ZCTA is one vertex.
Let $F$ be the scenario's fine channels, $S$ the states (including DC), and
$\mathcal H$ the permitted routing structures. For $h\in\mathcal H$,
$C_h$ is its set of planning channels and
$\rho_h:Z\times F\to C_h$ is its total routing map, obeying the settled
national/fallback rules. The permitted structures are an input, not
necessarily the five structures previously tested; $|\mathcal H|=1$
means a fixed routing structure. Set
\begin{equation}\label{eq:data}
\begin{aligned}
 Z_{hc}&=\{z\in Z:\exists f\in F,\ \rho_h(z,f)=c\},\\
 d_{hcz}&=\sum_{f:\rho_h(z,f)=c} r_f m_{zf},\qquad
 Z^0_{hc}=\{z\in Z_{hc}:d_{hcz}=0\}.
\end{aligned}
\end{equation}
Here $m_{zf}\ge0$ is descaled opportunity and $r_f$ is its dollar rate from
\emph{the whole extract}, never recomputed after regional extraction.
All monetary quantities below, including $r_f m_{zf}$, are in millions of dollars.
Further data are land area $a_z\ge0$, state $s(z)$, the P-adjacency graph
$G=(Z,E_P)$, and the width graph data described in Section~\ref{sec:neck}.
$E_P$ contains polygon rook edges (positive shared boundary, not a corner),
approved connectors, and committed R5 land-only gap adjacency. R5 joins
same-state ZCTAs facing each other across land in no ZCTA, never across water.
The land-only adjacency file is not yet shipped; see \MODEL, \S9.

Let $\mathcal M\subseteq Z$ be the Manhattan ZCTAs: rows of
\path{reference/2025/zcta_reference.csv.gz} whose primary county is
\texttt{36061} (New York County), intersected with $Z$. This is the
2025 reference-table county assignment, not a ZIP-prefix rule.
The primary-county convention is documented at \path{td/geo.py:45--46}:
\begin{quote}\scriptsize
\verb|urban-area overlays run on those land pieces.  A ZCTA's primary county, place and|\\
\verb|urban area are the ones holding the most of its land; its state is its primary county's, and|
\end{quote}
The numerical map evaluators $\Phi_{\rm shape}(\lambda)$ and
$\Phi_{\rm balance}(\lambda)$ are \textbf{owner inputs not yet defined}.
No surrogate evaluator or calibrated shape gate is silently substituted.

\subsection*{Notation}
\begin{center}
\begin{tabular}{p{0.27\textwidth}p{0.65\textwidth}}
\hline
Symbol & Meaning\\ \hline
$Z,F,S;\ \mathcal H,C_h$ & ZCTAs, fine channels, states; permitted structures and planning channels.\\
$Z_{hc},Z^0_{hc};\ \mathcal M$ & Channel footprint, its zero-dollar ZCTAs; Manhattan ZCTAs.\\
$d_{hcz},a_z,s(z)$ & Routed dollars, land area, state.\\
$J=\{1,\ldots,54\}$ & District labels available in each channel.\\
$q_h,y_{hcj},x_{hcjz}$ & Choose structure, activate district, own ZCTA (binary).\\
$b_{hcjs},w_{hcj},p_{hcs}$ & Touch state, WIFI allowance, split state (binary).\\
$r_{hcjz},t_{hcjz};\ f_{hcj,uv}$ & Ordinary root, optional Manhattan root (binary); nonnegative directed flow.\\
$X_{hcj},D_{hcj},K_{hc}$ & Owned set, drawn dollars, active district count.\\
$H_{hc}(X),R_{hc}(X)$ & Width graph and its outside zero-opportunity relays.\\ \hline
\end{tabular}
\end{center}
Use $t$ for the Manhattan root to distinguish it from state $s(z)$;
$r_{hcjz}$ is also distinct from rate $r_f$.
All district-indexed rows apply to $h\in\mathcal H$, $c\in C_h$, $j\in J$;
ZCTA rows use $z\in Z_{hc}$ and state rows use $s\in S$ unless stated otherwise.
Write $X_{hcj}=\{z:x_{hcjz}=1\}$ and $a(A)=\sum_{z\in A}a_z$.

\section{Routing, coverage and district count}
\begin{equation}\label{eq:coverage}
\begin{aligned}
 \sum_{h\in\mathcal H}q_h&=1,& y_{hcj}&\le q_h,\\
 \sum_{j\in J}x_{hcjz}&=q_h,& x_{hcjz}&\le y_{hcj},\\
 \sum_{z\in Z_{hc}}x_{hcjz}&\ge y_{hcj},&
 K_{hc}:=\sum_{j\in J}y_{hcj}&\ge q_h,\\
 48&\le\sum_{h\in\mathcal H}\sum_{c\in C_h}K_{hc}\le54.
\end{aligned}
\end{equation}
For the chosen structure $h$, the routed ledger is
\begin{equation}\label{eq:ledger}
 \lambda(z,f)=(h,\rho_h(z,f),j)
 \quad\text{where }x_{h,\rho_h(z,f),j,z}=1.
\end{equation}
Thus fine channels routed together have the same ZCTA owner. Every cell of
$Z\times F$ is assigned exactly once, including zero-opportunity cells;
no cell is inferred from polygon availability. Every selected planning
channel has at least one nonempty district. F0 has no dropped-channel exception.

\section{Hard dollar windows and the conditional WIFI cap}
Define drawn dollars and window data by
\begin{equation}\label{eq:dollars}
 D_{hcj}=\sum_{z\in Z_{hc}}d_{hcz}x_{hcjz}.
\end{equation}
\begin{center}
\begin{tabular}{lrrl}
\hline
Channel & $L_c$ & $U_c$ & Conditional cap\\ \hline
national & 1000 & 1437.5 & none\\
WH & 850 & 1150 & none\\
FI & 765 & 1035 & none\\
WIFI & 850 & 1150 & 1600, whole-state districts only\\ \hline
\end{tabular}
\end{center}
The hard bounds are
\begin{equation}\label{eq:windows}
\begin{aligned}
 L_cy_{hcj}\le D_{hcj}&\le U_cy_{hcj} &&(c\ne\mathrm{WIFI}),\\
 850y_{hcj}\le D_{hcj}&\le1150y_{hcj}+450w_{hcj} &&(c=\mathrm{WIFI}).
\end{aligned}
\end{equation}
Exact state incidence and allowance eligibility are
\begin{equation}\label{eq:whole}
\begin{aligned}
 x_{hcjz}&\le b_{hcj,s(z)},&
 b_{hcjs}&\le\sum_{z\in Z_{hc}:s(z)=s}x_{hcjz},\\
 w_{hcj}&\le y_{hcj},& w_{hcj}&=0 &&(c\ne\mathrm{WIFI}),\\
 x_{hcjz}&\ge w_{hcj}+b_{hcj,s(z)}-1 &&&&(c=\mathrm{WIFI}).
\end{aligned}
\end{equation}
If $w=1$, all footprint ZCTAs in every touched state belong to that district.
By coverage, this is exclusive ownership of each touched state in that channel;
several whole states may form one district. The allowance is optional and
never grants 1600 to a split-state WIFI district. Zero-opportunity ZCTAs count.

F0 uses exact dollars. The production scorer instead rounds to whole dollars;
this is a declared numerical convention, not a solver tolerance:
\path{tools/looks/score.py:337}
\begin{quote}\small
\verb|inside = lo is not None and round(lo) <= round(x) <= round(hi)|
\end{quote}
The cap and ownership test are at \path{tools/looks/score.py:111,334}:
\begin{quote}\fontsize{7}{9}\selectfont
\verb|WHOLE_STATE_CAP = {"WIFI": 1.6e9}       # a district of whole states only (owner, 2026-10-13)|\\
\verb|whole = bool(st) and all(owners[ch, s] == {d} for s in st)|
\end{quote}

\section{P-connectivity and the Manhattan detached-piece exception}\label{sec:P}
For each channel footprint set $n=|Z_{hc}|$. Use both directed arcs of every
edge of $G[Z_{hc}]$. Impose one ordinary root and at most one additional root,
which must be an owned Manhattan ZCTA:
\begin{equation}\label{eq:roots}
\begin{aligned}
 \sum_{z\in Z_{hc}}r_{hcjz}&=y_{hcj},& r_{hcjz}&\le x_{hcjz},\\
 \sum_{z\in Z_{hc}}t_{hcjz}&\le y_{hcj},&t_{hcjz}&\le x_{hcjz},\\
 t_{hcjz}&=0 &&(z\notin\mathcal M).
\end{aligned}
\end{equation}
\begin{equation}\label{eq:flow}
\begin{aligned}
 0\le f_{hcj,uv}&\le(n-1)x_{hcju},&
 f_{hcj,uv}&\le(n-1)x_{hcjv},\\
 \sum_{u:uz\in E_P}f_{hcj,uz}-\sum_{v:zv\in E_P}f_{hcj,zv}
 &\ge x_{hcjz}-n(r_{hcjz}+t_{hcjz}).
\end{aligned}
\end{equation}
The sums in \eqref{eq:flow} use only footprint vertices. These rows require
one connected district, except that it may have one additional component
containing Manhattan. They do not waive a third component or any other
condition merely because a district touches Manhattan. This implements the
owner's 15 October ruling: ``for manhattan a detached piece that contains
manhattan is fine'' (\issue{140}).

\textbf{Proof of equivalence.} Summing flow balance over a selected component
without either root gives $0\ge|A|$, a contradiction. There are at most two
roots, and the optional root is in Manhattan. Conversely, for a connected
nonempty set choose one ordinary root and set $t=0$. For two components,
one containing Manhattan, put the optional root in that component and the
ordinary root in the other. Route one unit along a rooted spanning tree to
each nonroot vertex in each component. No edge carries more than $n-1$,
and each root can supply its component's demand. Empty districts have no
roots or flow. Thus when $X\cap\mathcal M=\varnothing$ this is exactly the
ordinary single-root formulation. The statement is about graph components;
a multipart ZCTA remains one vertex.

Zero-opportunity relays do \emph{not} create P-paths: a zero-opportunity ZCTA
must be owned by the district to join its pieces. Relays enter only the
width graph below. Approved connectors and committed R5 land-only edges
are genuine P-edges, unlike width-only gap edges.

\section{N-connectivity: widths, relays and a finite constraint family}\label{sec:neck}
For a candidate owned set $X\subseteq Z_{hc}$, let
$R_{hc}(X)\subseteq Z^0_{hc}\setminus X$ consist of zero-opportunity ZCTAs
reachable from $X$ along chains of zero-opportunity ZCTAs in the width
adjacency data. They may be owned by another district. Construct
$H_{hc}(X)$ on $X\cup R_{hc}(X)$, using these rules:
\begin{enumerate}
\item Polygon shared borders and committed gap widths contribute passage
width; finite widths are floored to whole centimetres.
\item Every approved connector has unlimited width, always. The older
``land would do'' zero-width rule is not part of F0.
\item Land in no ZCTA contributes gap width. For reproduction of the current
width data, \MODEL, \S9 R2 uses land-only gap rows when present, otherwise
legacy Voronoi dropped-edge widths, which lack a land/water marker.
Those legacy edges widen necks only; they never join P-components.
\item Outside relays have zero land and mass. A cut may place each relay on
either side; a route through relays is no wider than its narrowest cut.
Positive-opportunity ZCTAs outside $X$ are not relays.
\end{enumerate}
For $U\subseteq V(H)$, $\delta_H(U)$ is the set of edges with one endpoint
in $U$ and the other outside; $\capcut_H(\delta_H(U))$ is their total width.
For every nonempty proper $A\subset X$ connected on $G[A]$, require
\begin{equation}\label{eq:neck}
\begin{gathered}
 0.05a(X)\le a(A)\le0.5a(X),\qquad A\cap\mathcal M=\varnothing\\
 \Longrightarrow\quad
 \min_{T\subseteq R_{hc}(X)}
 \capcut_{H_{hc}(X)}\bigl(\delta_{H_{hc}(X)}(A\cup T)\bigr)\ge W,\\
 W=1{,}000{,}000\ \mathrm{cm}=10\ \mathrm{km}.
\end{gathered}
\end{equation}
The cut-off part is one connected smaller side holding at least 5\% and
at most 50\% of the district's land. Its complement may be disconnected.
Mass is reported, not used to decide a neck. Distinct small fringes joined
only through relays do not combine into one qualifying part.

\textbf{Manhattan neck interpretation.} In \eqref{eq:neck}, the Manhattan
neck is identified by its smaller connected side containing a Manhattan
ZCTA. This is the supervisor's explicit reading of the IFA precedent, to
be confirmed with the owner, not an additional quotation from the owner.
Only that side is exempt: every qualifying side disjoint from Manhattan
still must pass, even in a district containing Manhattan. The detached-piece
exception in Section~\ref{sec:P} and this neck exception do not waive
coverage, windows, or any other failure.

An equivalent finite algebraic encoding is obtained by predefining
$\mathcal B_{hc}$ as all nonempty footprint subsets with a nonexempt
violation of \eqref{eq:neck}. For every $B\in\mathcal B_{hc}$ and $j\in J$,
impose
\begin{equation}\label{eq:nogood}
 \sum_{z\in B}(1-x_{hcjz})+
 \sum_{z\in Z_{hc}\setminus B}x_{hcjz}\ge y_{hcj}.
\end{equation}
\textbf{Proof.} For an active district the left side is the Hamming distance
between $X_{hcj}$ and $B$, so the row excludes exactly $X_{hcj}=B$.
Enumerating all and only the nonexempt forbidden sets encodes exactly
\eqref{eq:neck}. This is finite, not a tractability claim. A time-limited
neck search is \emph{unknown}, neither a mathematical violation nor a
feasibility proof. The production raw gate has no Manhattan-specific
acceptance path; an owner-applied exception must be distinguished from
its raw failure (\MODEL, \S9).

\section{Splits and the lexicographic objective}
Define state incidence counts and impose
\begin{equation}\label{eq:splits}
 R_{hcs}:=\sum_{j\in J}b_{hcjs},\qquad
 2p_{hcs}\le R_{hcs}\le q_h+53p_{hcs}.
\end{equation}
These inequalities make $p_{hcs}=1$ exactly when at least two districts
own territory in that channel-state, including zero-opportunity ZCTAs.
A state split in two planning channels counts twice, while five holders
in one channel still count as one split. Unselected structures have no
owners and no split indicators.

The reference objective is
\begin{equation}\label{eq:objective}
 \lexmin\left(
 \Phi_{\rm shape}(\lambda),\quad
 \Phi_{\rm balance}(\lambda),\quad
 \sum_{h\in\mathcal H}\sum_{c\in C_h}\sum_{s\in S}p_{hcs}
 \right).
\end{equation}
M1 (with only its stated exceptions) and the windows are hard constraints.
Shape has no numerical hard gate in this F0 statement; its evaluator and
the balance evaluator await owner definition. There is no automatic
$\eta$, free list, support-size/diameter cap, regional seal, rule C, or split
cap. K moves within its range, not by relaxing the window. With routing
and K fixed, feasibility separates by channel; a map-level objective
separates only if its evaluators do. Lexicographic minimization fixes each
previous optimum before minimizing the next component; optimizing only
shape does not certify the full tuple.

\section{Implementation boundary, pipeline and rule sources}
\textbf{Relation to the pipeline.} \MODEL\ describes F2--F142: support
master and drawability, drawing/repair, transport, audit, glue/sealing,
and branch-and-cut. Its stage contracts identify restrictions,
relaxations, or heuristics relative to F0 and state each certificate's
domain. A combination of restrictions and relaxations need not have
feasible-set containment in either direction. A fixed-group optimum,
a unit-level floor, or a regional feasibility result is not a national
ZIP-level F0 bound or feasibility certificate.

\textbf{Code entry boundary.} No production function solves F0.
The raw M1 check starts at \path{td/audit.py:940}:
\begin{quote}\scriptsize
\verb!def check_m1(run: Run, gap_width: bool | None = None, land_rule: bool | None = None) -> Check:!
\end{quote}
Routing is supplied by the source audit, not established by the run-folder
gate alone. Ranking is not the F0 objective:
\path{tools/looks/score.py:549} is \verb|def rank_key(s: dict) -> tuple:|;
line 550 begins with splits, not shape:
\begin{quote}\fontsize{7}{9}\selectfont
\verb|return (s["splits"], s["defects"], s["largest_extent_km"], s["states_per_district"], s["worst_dev"],|
\end{quote}
These differences between implementation and reference do not change F0
or the owner-owned \texttt{PROBLEM.md}.

\textbf{Executable specification.} \path{td/formulation/f0.py:18}:
\begin{quote}\scriptsize
\verb|def build_model(*, Z, F, rho, m, rates, state, E_P, manhattan, forbidden, shape, balance):|
\end{quote}
The Pyomo specification implements one supplied routing structure,
$J=1,\ldots,54$, and total $K\in[48,54]$, for at most 300 ZCTAs.
The caller supplies $\mathcal M$, the complete nonexempt family
$\mathcal B_{hc}$, and symbolic shape/balance expressions. Geometry and
owner numerical objectives are not invented by the builder.
WIFI uses the genuine GDP disjunction $w=0$ or
$(w=1\ \land\ x_{hcjz}\ge b_{hcj,s(z)}\ \forall z)$; other rows are algebraic.
Three ordered objectives are retained; only shape is initially active.
The small-instance tests exercise coverage, 48 active districts, dollars,
the conditional cap, forbidden sets, and ordinary/Manhattan-root flow.
They do not claim conformance with an unavailable production F0 optimizer;
the production component decomposition is checked separately from the
owner exception applied to it.

\subsection*{Source of each rule}
The following records were read for this transcription; quoted owner words
above are limited to the recorded Manhattan ruling. The table of numerical
windows and total K are retained from the approved \MODEL, former \S0,
not inferred from every historical grid-specific range in these issues.
\begin{enumerate}
\item \issue{57}: M1's authoritative 2025 polygon rook graph. Coverage,
multipart vertices, neck parameters, and gap width are recorded in the
M1 row of \texttt{PROBLEM.md}, citing \issue{108}, \issue{114},
\issue{121}, and \issue{131}.
\item \issue{112}: the band stays hard and K moves; direct dollar-window
rows replace a stricter symmetric surrogate. Numerical windows and
WIFI's conditional allowance match \path{tools/looks/score.py:109--111}
and ownership at line 334, subject to the rounding distinction above.
\item \issue{134}: the recorded priority is contiguity, shape, balance,
splits. This yields hard contiguity/windows and the lexicographic
objective, not the production splits-first ranking.
\item \issue{140}: approved connectors always have unlimited width;
zero-opportunity relays widen neck passages but do not transfer ownership;
land-only missing-land edges join pieces; Manhattan alone is excepted,
including its detached piece by the 15 October ruling.
\item \issue{135}: shape/arm calibration remains an owner-input boundary.
The record calls for an arm-share limit whose value must be calibrated
and confirmed. F0 does not invent that value or either numerical evaluator.
\end{enumerate}
\end{document}
